% GENERATED FILE. Edit canonical lessons, then run npm run generate:books.
% canonical-source-hash: fnv1a64:6371c22f1b6ea118
% canonical-lessons: FR-C413-voie, FR-C413-composter, FR-C413-controleur, FR-C413-wagon, FR-C413-desservir

\chapter{Sur le quai}
\label{ch:fr-sur-le-quai}
\hlchaptermodality{413}

\begin{chapteropening}
\textbf{By the end of this chapter:} \emph{I can find my train: the track, stamping a ticket, the ticket inspector, the carriage and the places a train serves.}

Complete the last lesson of chapter 413: I can find my train: the track, stamping a ticket, the ticket inspector, the carriage and the places a train serves.
\end{chapteropening}

\section[la voie]{la voie — the track, the platform number; the way}

\label{lesson:FR-C413-voie}



\begin{quote}
Before the new one: say the French for the flat-rate pass, then the French for to validate.
\end{quote}

\subsection*{You'll want to know: la voie}
\textbf{la voie} — \textquotedblleft{}the track, the platform number; the way\textquotedblright{}.

\begin{quote}
Le train de sept heures part de la voie six.
\end{quote}

\begin{quote}
\emph{The seven o'clock train leaves from platform six.}
\end{quote}

\emph{La voie} is the track with its number; \emph{le quai} is where you stand.

\begin{cousinweb}
\emph{voie} is inherited through Old French \emph{voie} from Latin \emph{via} 'road, way' (EtymDB).
\end{cousinweb}

\begin{grammarlens}[title={What you've built}]
One more word for telling, explaining and describing at B1.
\end{grammarlens}

\subsection*{Your turn}
\begin{itemize}
\raggedright
  \item \emph{Say it:} \emph{la voie}
  \item \emph{Say it:} \emph{la voie}, once more
  \item \emph{Say it:} a sentence of your own with \emph{la voie}
  \item \emph{Recall:} say the French for the flat-rate pass, then the French for to validate, then say \emph{la voie} again
\end{itemize}

\begin{tcolorbox}[breakable,colback=teal!4,colframe=teal!35!black,title={Before you move on}]
What does la voie mean? (\textquotedblleft{}The track, the platform number; the way\textquotedblright{}.) Say it once more.
\end{tcolorbox}
\section[composter]{composter — to stamp (a ticket)}

\label{lesson:FR-C413-composter}



\begin{quote}
Before the new one: say the French for to validate, then the French for the track.
\end{quote}

\subsection*{You'll want to know: composter}
\textbf{composter} — \textquotedblleft{}to stamp (a ticket)\textquotedblright{}.

\begin{quote}
Avant de monter, il a composté son billet.
\end{quote}

\begin{quote}
\emph{Before getting on, he stamped his ticket.}
\end{quote}

Paper tickets were stamped in a small machine on the platform; in Switzerland, \emph{oblitérer} is also heard.

\begin{cousinweb}
\emph{composter} is built on \emph{composteur}, a stamping machine, which is usually traced to Italian \emph{compositore} 'typesetter', from Latin \emph{compōnere} 'to put together'. EtymDB lists the verb without a source.
\end{cousinweb}

\begin{grammarlens}[title={What you've built}]
One more word for telling, explaining and describing at B1.
\end{grammarlens}

\subsection*{Your turn}
\begin{itemize}
\raggedright
  \item \emph{Say it:} \emph{composter}
  \item \emph{Say it:} \emph{composter}, once more
  \item \emph{Say it:} a sentence of your own with \emph{composter}
  \item \emph{Recall:} say the French for to validate, then the French for the track, then say \emph{composter} again
\end{itemize}

\begin{tcolorbox}[breakable,colback=teal!4,colframe=teal!35!black,title={Before you move on}]
What does composter mean? (\textquotedblleft{}To stamp (a ticket)\textquotedblright{}.) Say it once more.
\end{tcolorbox}
\section[le contrôleur]{le contrôleur — the ticket inspector}

\label{lesson:FR-C413-controleur}



\begin{quote}
Before the new one: say the French for the track, then the French for to stamp.
\end{quote}

\subsection*{You'll want to know: le contrôleur}
\textbf{le contrôleur} — \textquotedblleft{}the ticket inspector\textquotedblright{}.

\begin{quote}
Le contrôleur passe dans le train et vérifie les billets.
\end{quote}

\begin{quote}
\emph{The inspector goes through the train and checks the tickets.}
\end{quote}

Feminine: \emph{la contrôleuse}.

\begin{cousinweb}
\emph{contrôleur} comes from Old French \emph{contrerollour} (EtymDB), built on \emph{contrerole} 'a second register for checking', from Medieval Latin \emph{contrārotulum}.
\end{cousinweb}

\begin{grammarlens}[title={What you've built}]
One more word for telling, explaining and describing at B1.
\end{grammarlens}

\subsection*{Your turn}
\begin{itemize}
\raggedright
  \item \emph{Say it:} \emph{le contrôleur}
  \item \emph{Say it:} \emph{le contrôleur}, once more
  \item \emph{Say it:} a sentence of your own with \emph{le contrôleur}
  \item \emph{Recall:} say the French for the track, then the French for to stamp, then say \emph{le contrôleur} again
\end{itemize}

\begin{tcolorbox}[breakable,colback=teal!4,colframe=teal!35!black,title={Before you move on}]
What does le contrôleur mean? (\textquotedblleft{}The ticket inspector\textquotedblright{}.) Say it once more.
\end{tcolorbox}
\section[le wagon]{le wagon — the carriage (of a train)}

\label{lesson:FR-C413-wagon}



\begin{quote}
Before the new one: say the French for to stamp, then the French for the ticket inspector.
\end{quote}

\subsection*{You'll want to know: le wagon}
\textbf{le wagon} — \textquotedblleft{}the carriage (of a train)\textquotedblright{}.

\begin{quote}
Notre place est dans le wagon numéro douze.
\end{quote}

\begin{quote}
\emph{Our seat is in carriage number twelve.}
\end{quote}

Rail staff often say \emph{la voiture}; most people say \emph{le wagon}. \emph{Le wagon-restaurant} — the dining car.

\begin{cousinweb}
\emph{wagon} was borrowed from English \emph{wagon}, which came from Dutch \emph{wagen} (EtymDB).
\end{cousinweb}

\begin{grammarlens}[title={What you've built}]
One more word for telling, explaining and describing at B1.
\end{grammarlens}

\subsection*{Your turn}
\begin{itemize}
\raggedright
  \item \emph{Say it:} \emph{le wagon}
  \item \emph{Say it:} \emph{le wagon}, once more
  \item \emph{Say it:} a sentence of your own with \emph{le wagon}
  \item \emph{Recall:} say the French for to stamp, then the French for the ticket inspector, then say \emph{le wagon} again
\end{itemize}

\begin{tcolorbox}[breakable,colback=teal!4,colframe=teal!35!black,title={Before you move on}]
What does le wagon mean? (\textquotedblleft{}The carriage (of a train)\textquotedblright{}.) Say it once more.
\end{tcolorbox}
\section[desservir]{desservir — to stop at, to serve (a place); to clear away (after a meal)}

\label{lesson:FR-C413-desservir}



\begin{quote}
Before the new one: say the French for the ticket inspector, then the French for the carriage.
\end{quote}

\subsection*{You'll want to know: desservir}
\textbf{desservir} — \textquotedblleft{}to stop at, to serve (a place); to clear away (after a meal)\textquotedblright{}.

\begin{quote}
Ce train dessert tous les villages de la vallée.
\end{quote}

\begin{quote}
\emph{This train serves all the villages in the valley.}
\end{quote}

Also \emph{desservir} after a meal — to clear away the dishes.

\begin{cousinweb}
\emph{desservir} is built from \emph{des-} and \emph{servir} 'to serve', which is inherited from Latin \emph{servīre} (EtymDB).
\end{cousinweb}

\begin{grammarlens}[title={What you've built}]
One more word for telling, explaining and describing at B1.
\end{grammarlens}

\subsection*{Your turn}
\begin{itemize}
\raggedright
  \item \emph{Say it:} \emph{desservir}
  \item \emph{Say it:} \emph{desservir}, once more
  \item \emph{Say it:} a sentence of your own with \emph{desservir}
  \item \emph{Recall:} say the French for the ticket inspector, then the French for the carriage, then say \emph{desservir} again
\end{itemize}

\begin{tcolorbox}[breakable,colback=teal!4,colframe=teal!35!black,title={Before you move on}]
What does desservir mean? (\textquotedblleft{}To stop at, to serve (a place); to clear away (after a meal)\textquotedblright{}.) Say it once more.
\end{tcolorbox}
